\begin{table*}[t]
\caption{ST-MTM geometry metrics on 3-D inputs. Lower SNS and TD and higher TW are better.}\label{tab:geometry3d}
\centering\scriptsize\setlength{\tabcolsep}{3pt}
\begin{tabular}{lrrrrrrrrrrrr}\toprule
 & \multicolumn{3}{c}{TMTM} & \multicolumn{3}{c}{ST-MTM} & \multicolumn{3}{c}{RG-MTM-P} & \multicolumn{3}{c}{RG-MTM-R20}\\
Data & SNS & TW$_1$ & TD & SNS & TW$_1$ & TD & SNS & TW$_1$ & TD & SNS & TW$_1$ & TD\\\midrule
Tangaroa & 0.5255 & 0.6701 & 64.8553 & 0.1655 & 0.8137 & 138.5519 & 0.3571 & 0.7271 & 158.2670 & 0.3353 & 0.7458 & 39.7093\\
Cylinder & 0.2956 & 0.8017 & 85.0339 & 0.1346 & 0.8989 & 43.1663 & 0.1393 & 0.7838 & 72.3427 & 0.0399 & 0.8912 & 29.1130\\
\bottomrule\end{tabular}
\par\smallskip{\footnotesize\raggedright SNS uses source support-centroid distances; TW uses the common $k=1$ and only frames with more than two features. TD sums matched displacements in full-canvas units, with one fixed mapping per sequence. -- denotes static data with no temporal TD; Benzene section displacement is reported separately in the supplement. Eligible frame counts and complete precision are retained in Supplement S3.\par}
\end{table*}
